%%%%%%%%%%%%%%%%%%%%% Chapter 21 %%%%%%%%%%%%%%%%%%%%%%%%%%%%%%%%%
% Zero-Knowledge Machine Learning (ZK-ML)
%%%%%%%%%%%%%%%%%%%%%%%%%%%%%%%%%%%%%%%%%%%%%%%%%%%%%%%%%%%%%%%%%
\motto{Zero-Knowledge Machine Learning (ZK-ML) allows you to prove that a model was executed correctly without ever revealing the model's weights or the user's private inputs.}

\chapter{Zero-Knowledge Machine Learning (ZK-ML)}
\label{ch:zk-ml}

\abstract{Building on the foundations of privacy and security, this chapter introduces the third pillar of trust: Verifiability. Zero-Knowledge Machine Learning (ZK-ML) uses sophisticated cryptographic proofs to prove that a specific machine learning model was executed correctly on specific input data, producing a specific output, all while keeping the inputs and model weights completely private. We analyze the transition of neural networks into arithmetic circuits, the Rank-1 Constraint System (R1CS), and the practical implementation of verifiable AI using the ezkl framework. By enabling high-performance server-side execution with lightweight client-side verification, ZK-ML provides the ultimate architectural tool for accountability in the age of AI.}

\section{The Verifiability Gap}
\label{sec:verifiability_gap_ch19}

In previous chapters, we have seen how to protect data (Differential Privacy) and how to training models securely (Federated Learning). However, a fundamental gap remains: How can a client verify that the model they are querying is actually the one they were promised? How can an auditor verify that a model is unbiased without requiring the company to reveal its proprietary weights or its sensitive training data?

Traditional encryption protects data at rest and in transit, but it does not prove that a computation happened correctly. \textbf{Zero-Knowledge Machine Learning (ZK-ML)} \cite{shrestha2023} solves this by generating a cryptographic proof that a computation was performed according to a predefined "circuit" of operations.

\section{How ZK-ML Works: From Neural Networks to Arithmetic Circuits}
\label{sec:zkml_mechanics_ch19}

To prove a model's execution, ZK-ML converts the neural network's layers into an arithmetic circuit that can be validated using \textbf{zk-SNARKs} (Zero-Knowledge Succinct Non-Interactive Arguments of Knowledge). The computation is expressed as a \textbf{Rank-1 Constraint System (R1CS)}, where each operation must satisfy a specific polynomial constraint:
\begin{equation}
\langle a_i, w \rangle \cdot \langle b_i, w \rangle = \langle c_i, w \rangle
\end{equation}
Here, $w$ represents the "witness"---the collection of all inputs, weights, and intermediate values. The Prover uses their secret witness to generate a proof $\pi$, while the Verifier uses a public verification key to check that the constraints were satisfied.

\section{The Power of Asymmetry}
\label{sec:asymmetry_ch19}

One of the most powerful properties of ZK-ML is its asymmetry. While generating the proof (the "Prover" side) is computationally intensive and often requires server-class hardware, verifying the proof (the "Verifier" side) takes only milliseconds and consumes negligible memory. This allows lightweight edge devices, such as smartphones or IoT sensors, to verify that a massive server-side model has processed their data correctly without ever seeing the model's internals.

\section{Applications of Verifiable AI}
\label{sec:zkml_applications_ch19}

\begin{table}
\caption{ZK-ML Architectural Benefits}
\label{tab:zkml_benefits_ch19}
\begin{tabular}{p{3.5cm}p{7cm}}
\hline\noalign{\smallskip}
Use Case & Architect's Benefit \\
\noalign{\smallskip}\svhline\noalign{\smallskip}
Private Inference & User data stays encrypted; server proves correct computation. \\
Verifiable Ownership & Prove model integrity and pedigree without exposing IP weights. \\
Compliance & Prove model fairness to auditors without exposing dataset details. \\
\noalign{\smallskip}\hline\noalign{\smallskip}
\end{tabular}
\end{table}

\section{Conclusion}
\label{sec:conclusion_zkml_ch19}

Zero-Knowledge Machine Learning transforms AI from a "trust-me" black box into a verifiably correct utility. By anchoring trust in cryptographic proof rather than institutional promises, we enable a new class of secure and accountable applications. While ZK-ML provides a proof of \textit{execution}, it does not necessarily prove the model's \textit{correctness} or \textit{safety} properties in all possible states. To achieve that, we must turn to \textbf{Formal Verification}, which we explore in the next chapter.


